% Frahan StonePack — algorithm math derivations, structured for later Lean formalization.
% Each algorithm: formal Definitions, then Theorems/Lemmas with proofs or proof-sketches.
% The Definition/Theorem/Proof structure maps onto Lean `def` / `theorem` / proof terms.
% Compile: pdflatex frahan_algorithm_derivations.tex
\documentclass[11pt]{article}
\usepackage{amsmath,amssymb,amsthm,mathtools,enumitem,geometry,hyperref}
\geometry{margin=1in}
\theoremstyle{definition}\newtheorem{definition}{Definition}[section]
\theoremstyle{plain}\newtheorem{theorem}{Theorem}[section]
\newtheorem{lemma}[theorem]{Lemma}\newtheorem{proposition}[theorem]{Proposition}
\newtheorem{corollary}[theorem]{Corollary}
\newcommand{\R}{\mathbb{R}}\newcommand{\Z}{\mathbb{Z}}\newcommand{\Norm}[1]{\lVert #1\rVert}
\newcommand{\inner}[2]{\langle #1,#2\rangle}\DeclareMathOperator{\area}{area}
\DeclareMathOperator{\conv}{conv}\DeclareMathOperator*{\argmin}{arg\,min}
\title{Frahan StonePack — Algorithm Derivations\\ \large (LaTeX source for Lean proof management)}
\author{Frahan StonePack}\date{2026-06-18}
\begin{document}\maketitle
\noindent\textbf{Scope.} Formal math for the algorithms implemented in the 2026-06-18 session:
GPR fracture-surface kriging (regression kriging, $k$-planes, robust clip, hull mask), Panel Tile
Surface (best-fit-plane planarization), and Slab Trim (convex greedy trim, Imai--Iri concave kerf-follow)
plus concave nesting. Definitions and theorems are written so each maps to a Lean \texttt{def}/\texttt{theorem}.

\section{Regression kriging of a fracture bed}
\begin{definition}[Bed data]
A bed is a finite set of picks $\{(x_i,y_i,d_i)\}_{i=1}^n$, $(x_i,y_i)\in\R^2$ planimetric, $d_i\in\R$ depth.
Write $p_i=(x_i,y_i)$.
\end{definition}
\begin{definition}[Least-squares dip plane]
The dip plane is $\pi(p)=a^\top p + c$ with $(a,c)=(a_1,a_2,c)\in\R^3$ minimizing
$\sum_i (d_i-a^\top p_i-c)^2$. Equivalently $(a,c)$ solves the normal equations $M\theta=b$ where
$M=\sum_i \tilde p_i\tilde p_i^\top$, $b=\sum_i d_i\tilde p_i$, $\tilde p_i=(x_i,y_i,1)$.
\end{definition}
\begin{lemma}[Existence/uniqueness]\label{lem:plane}
If the $p_i$ are not all collinear then $M\succ0$ and the dip plane is unique.
\end{lemma}
\begin{proof}
$M=\sum_i\tilde p_i\tilde p_i^\top\succeq0$. For $u\in\R^3$, $u^\top M u=\sum_i(\tilde p_i^\top u)^2=0$
iff every $p_i$ lies on the line $u_1x+u_2y+u_3=0$; non-collinearity forbids this for $u\neq0$, so
$M\succ0$ and $\theta=M^{-1}b$ is unique.
\end{proof}
\begin{definition}[Residual and Gaussian covariance]
Residuals $r_i=d_i-\pi(p_i)$. With sill $\sigma^2=\mathrm{Var}(r)$, range $\rho>0$, nugget $\tau\ge0$, the
covariance is $C(h)=\sigma^2\exp(-h^2/\rho^2)+\tau\,\mathbf 1[h=0]$, and $K_{ij}=C(\Norm{p_i-p_j})$.
\end{definition}
\begin{definition}[Regression-kriging predictor]
$\hat d(q)=\pi(q)+k(q)^\top\alpha$ where $k(q)_i=\sigma^2\exp(-\Norm{q-p_i}^2/\rho^2)$ and $\alpha=K^{-1}r$.
(The residual term is the simple-kriging BLUP of $r$ at $q$.)
\end{definition}
\begin{theorem}[Exactness vs.\ smoothing]\label{thm:nugget}
If $\tau=0$ and the $p_i$ are distinct then $\hat d(p_j)=d_j$ for all $j$ (interpolation). If $\tau>0$ the
predictor is a strict smoother: the residual correction $\hat d(p_j)-\pi(p_j)=k(p_j)^\top(K_0+\tau I)^{-1}r$
has $\ell_2$ norm strictly less than $\Norm{r}$, decreasing in $\tau$.
\end{theorem}
\begin{proof}
$\tau=0$, distinct points $\Rightarrow K$ is the Gaussian kernel Gram matrix, SPD, and
$k(p_j)^\top K^{-1}=e_j^\top$ (row $j$ of $K K^{-1}$), so $\hat d(p_j)=\pi(p_j)+r_j=d_j$. For $\tau>0$,
$K=K_0+\tau I$ with $K_0\succeq0$; then $k(p_j)^\top(K_0+\tau I)^{-1}r$ has $\ell_2$-norm bounded by
$\Norm{r}\,\Norm{K_0}/(\lambda_{\min}(K_0)+\tau)<\Norm{r}$ scaling, monotonically decreasing in $\tau$.
\end{proof}
\noindent\emph{Remark (Lean).} Theorem~\ref{thm:nugget} is the formal statement of why a positive nugget
removes the Gaussian-GP oscillation while a near-zero nugget interpolates; it reduces to SPD linear algebra.

\section{$k$-planes bed separation}
\begin{definition}[$k$-planes assignment]
Given picks and $k$ planes $\pi_1,\dots,\pi_k$, the assignment is
$\ell(i)=\argmin_{c} |d_i-\pi_c(p_i)|$. A configuration $(\ell,\{\pi_c\})$ is \emph{stable} if each $\pi_c$
is the LS dip plane of $\{i:\ell(i)=c\}$ and no pick changes assignment.
\end{definition}
\begin{theorem}[Monotone convergence]\label{thm:kplanes}
Alternating (i) refit each $\pi_c$ from its members and (ii) reassign each $i$ to $\argmin_c|d_i-\pi_c(p_i)|$
does not increase the cost $J=\sum_i \min_c (d_i-\pi_c(p_i))^2$, and terminates at a stable configuration in
finitely many steps.
\end{theorem}
\begin{proof}
Step (i) minimizes $\sum_{i:\ell(i)=c}(d_i-\pi_c(p_i))^2$ over $\pi_c$ (Lemma~\ref{lem:plane}), so it cannot
raise $J$. Step (ii) replaces each term by a $\le$ term, so it cannot raise $J$. $J\ge0$ and there are
finitely many assignments $\ell$ ($k^n$); a strictly decreasing-or-equal sequence over a finite assignment
set with optimal planes per assignment cannot cycle, hence terminates. This is Lloyd's argument with
points replaced by affine 1-codimensional fits.
\end{proof}

\section{Robust $2\sigma$ clip}
\begin{definition}
Given residuals $r_i=d_i-\pi(p_i)$ with mean $\mu$, std $s$, the $2\sigma$-clip keeps
$I'=\{i: |r_i-\mu|\le 2s\}$, refits $\pi$ on $I'$, and repeats.
\end{definition}
\begin{proposition}
Each clip pass is non-expansive on the kept set ($I'\subseteq I$) and the per-bed residual variance is
non-increasing across passes until a fixed point ($I'=I$) is reached; with the $\ge4$-point floor it
terminates.
\end{proposition}

\section{Convex-hull surface mask}
\begin{definition}[Convex hull]
$\conv(S)$ for finite $S\subset\R^2$ is the smallest convex set containing $S$; its boundary vertices are
computed by Andrew's monotone chain in $O(n\log n)$.
\end{definition}
\begin{definition}[Hull mask]
A grid point $q$ is drawn iff $q\in\conv(\{p_i\})$ (expanded outward by margin $m$), tested by ray casting:
$q$ is inside iff a generic ray from $q$ crosses the boundary an odd number of times.
\end{definition}
\begin{proposition}[No extrapolation]
With the hull mask the reconstructed surface is supported on $\conv(\{p_i\})$; hence the bed is never
evaluated outside the planimetric convex hull of its own picks, eliminating extrapolation overshoot.
\end{proposition}

\section{Panel Tile Surface: planarization}
\begin{definition}[Quad and best-fit plane]
A panel quad is $Q=\{q_0,q_1,q_2,q_3\}\subset\R^3$. Its best-fit plane $\Pi$ has unit normal $\nu$ the
eigenvector of $\sum_k (q_k-\bar q)(q_k-\bar q)^\top$ for the least eigenvalue ($\bar q$ the centroid).
\end{definition}
\begin{definition}[Planarity and cut tile]
Planarity $\delta(Q)=\max_k \operatorname{dist}(q_k,\Pi)=\max_k|\nu^\top(q_k-\bar q)|$. The cut tile is the
projection $\hat q_k=q_k-(\nu^\top(q_k-\bar q))\nu$ mapped by the rigid motion $\Pi\!\to\!\{z=0\}$.
\end{definition}
\begin{lemma}[$\delta$ minimizes worst-case deviation among planes through $\bar q$]
$\Pi$ minimizes $\sum_k \operatorname{dist}(q_k,\cdot)^2$ over planes through $\bar q$; and $\delta(Q)$ is the
Chebyshev (max) deviation of the corners from that plane.
\end{lemma}
\begin{proof}
For a plane through $\bar q$ with normal $u$, $\sum_k\operatorname{dist}^2=u^\top\big(\sum_k(q_k-\bar q)(q_k-\bar q)^\top\big)u$,
minimized over $\Norm u=1$ at the least eigenvector (Rayleigh quotient). $\delta$ is then the max over $k$
of $|u^\top(q_k-\bar q)|$ by definition.
\end{proof}
\noindent\emph{Use.} $\delta(Q)$ is the fabrication facet error: a flat stone panel deviates from the
curved surface by at most $\delta$; refine $U\times V$ until $\delta\le$ tolerance.

\section{Slab Trim I --- convex blank (greedy half-plane)}
\begin{definition}[Polygon, signed area, reflex vertex]
A simple polygon $P=(v_0,\dots,v_{n-1})$ (indices mod $n$). Signed area
$A(P)=\tfrac12\sum_i (x_i y_{i+1}-x_{i+1}y_i)$; CCW iff $A>0$. Vertex $v_i$ is \emph{reflex} (CCW) iff
$(v_i-v_{i-1})\times(v_{i+1}-v_i)<0$.
\end{definition}
\begin{definition}[Half-plane clip, Sutherland--Hodgman]
For half-plane $H=\{x:\nu^\top(x-p)\ge0\}$, $\mathrm{clip}(P,H)$ outputs, per edge $(a,b)$: $a$ if inside;
and the segment--line intersection when the edge crosses $\partial H$.
\end{definition}
\begin{lemma}[Subset / material-removal]\label{lem:sh}
For any simple polygon $P$ and half-plane $H$, $\mathrm{clip}(P,H)=P\cap H\subseteq P$, and
$\area(\mathrm{clip}(P,H))\le\area(P)$. Thus every clip only removes material (a valid wire cut).
\end{lemma}
\begin{proof}
Sutherland--Hodgman against a single half-plane returns exactly the intersection of the (possibly
non-convex) subject polygon with the convex region $H$; intersection with $\R^2\supseteq H$ is a subset, and
area is monotone under $\subseteq$.
\end{proof}
\begin{definition}[Greedy convex trim]
While $P$ has a reflex vertex and cuts $<K$: pick the deepest reflex $v_i$; cut along the supporting line of
edge $(v_{i-1},v_i)$, keeping the half-plane containing the centroid; set $P\leftarrow\mathrm{clip}(P,H)$.
\end{definition}
\begin{theorem}[Convergence to a convex subset]\label{thm:trim}
The greedy convex trim produces a sequence $P=P_0\supseteq P_1\supseteq\cdots$ of polygons with
$A(P_{t+1})\le A(P_t)$ (Lemma~\ref{lem:sh}); on termination (no reflex vertex, or budget $K$) the output is a
convex polygon contained in $P$, hence the recovered blank is $P_{\mathrm{out}}=\bigcap_t H_t\cap P$ with
$\area(P_{\mathrm{out}})/\area(P)$ the recovered-area ratio (yield).
\end{theorem}
\begin{proof}
Each step clips by a half-plane through $v_{i-1}$ chosen to exclude the wedge that makes $v_i$ reflex, so the
turning at the affected boundary becomes left (convex); the number of reflex vertices is monotone
non-increasing and bounded below by $0$. If it reaches $0$ the polygon is convex (all cross products
$\ge0$). The output is an intersection of half-planes with $P$, which is convex, and $\subseteq P$ by
Lemma~\ref{lem:sh}. (Greedy is a heuristic for \emph{fewest} cuts; optimal-min-cut is \S7.)
\end{proof}

\section{Slab Trim II --- concave kerf-follow (Imai--Iri min-\#)}
\begin{definition}[$\varepsilon$-admissible chord]
For boundary vertices $v_i,v_j$ ($i<j$) the chord $\overline{v_iv_j}$ is $\varepsilon$-admissible iff every
intermediate vertex is within $\varepsilon$ of the chord: $\max_{i<k<j}\dfrac{|(v_k-v_i)\times(v_j-v_i)|}{\Norm{v_j-v_i}}\le\varepsilon$.
\end{definition}
\begin{definition}[Shortcut DAG and min-\# approximation]
$G=(V,E)$ with $V=\{0,\dots,n-1\}$ and $(i,j)\in E$ iff $\overline{v_iv_j}$ is $\varepsilon$-admissible. The
\emph{min-\# approximation} is a fewest-edges path $0\to n-1$ in $G$ (closed by the wrap edge), i.e.\ the
polyline with the fewest straight kerfs staying within $\varepsilon$ of the boundary.
\end{definition}
\begin{theorem}[Exactness of the DAG shortest path]\label{thm:imaiiri}
A breadth-first / fewest-edges shortest path in $G$ yields the minimum number of $\varepsilon$-admissible
straight segments approximating the chain; thus the kerf-follow cut count is \emph{optimal}, not merely
greedy. Complexity $O(n^2)$ time after $O(n^2)$ admissibility (improvable to $O(n^2)$ total / $O(n\log n)$
with cone tests).
\end{theorem}
\begin{proof}
Any valid $\varepsilon$-approximation is a chain $0=i_0<i_1<\cdots<i_m=n-1$ with each
$\overline{v_{i_t}v_{i_{t+1}}}$ admissible, i.e.\ a path in $G$ of $m$ edges; conversely any $G$-path is a
valid approximation. So valid approximations $\leftrightarrow$ $G$-paths, and ``fewest segments''
$\leftrightarrow$ ``fewest edges''. BFS computes a fewest-edges path exactly (every edge has unit weight),
giving the minimum. (This is the classical Imai--Iri min-\# result.)
\end{proof}
\begin{corollary}[Yield ordering]
With $\varepsilon$ small, kerf-follow recovers $\ge$ the convex trim's area (it keeps the concavities), at
the cost of $\ge$ as many cuts; the two modes bracket the cut-count/yield trade-off.
\end{corollary}

\section{Concave-in-concave nesting}
\begin{definition}[No-fit polygon, valid placement]
For a fixed sheet $S$ and a part $P$ under translation $t$ (and rotation $\theta$), the placement is valid iff
$R_\theta P + t \subseteq S$ and $(R_\theta P+t)$ is interior-disjoint from all previously placed parts. The
feasible $t$-set is the \emph{inner-fit region} $\mathrm{IFR}=\{t: R_\theta P+t\subseteq S\}$ minus the union
of no-fit polygons $\mathrm{NFP}(P_k,P)$ against placed parts.
\end{definition}
\begin{definition}[Raster heuristic]
Rasterize $S$ to a Boolean free-grid $F$ at cell $h$; rasterize $R_\theta P$ to a mask $M$ (dilated by the
kerf clearance). A bottom-left placement at offset $(o_x,o_y)$ is valid iff $M\wedge\lnot F[\text{window}]=\mathbf 0$;
place greedily largest-area-first, marking $F\leftarrow F\wedge\lnot M$.
\end{definition}
\begin{proposition}[Soundness and yield]
Every raster placement satisfies $R_\theta P+t\subseteq S$ and is interior-disjoint from placed parts up to
one cell ($h$); the fill ratio $\big(\sum_{\text{placed}}\area\big)/\area(S)$ lower-bounds the achievable
recovery and converges to the continuous NFP optimum as $h\to0$. Concavity of both $S$ and $P$ is handled
intrinsically (the masks encode arbitrary shape).
\end{proposition}

\section{2D nesting: no-fit polygon and bottom-left-fill}
\begin{definition}[Minkowski sum, no-fit polygon]
For $A,B\subset\R^2$, $A\oplus B=\{a+b:a\in A,b\in B\}$ and $-B=\{-b:b\in B\}$. The \emph{no-fit polygon}
of a fixed $A$ and a translating $B$ is $\mathrm{NFP}(A,B)=A\oplus(-B)$.
\end{definition}
\begin{theorem}[NFP non-overlap criterion]\label{thm:nfp}
Let $A,B$ be polygons with reference point of $B$ at $t$. Then $\operatorname{int}(A)\cap\operatorname{int}(B+t)=\varnothing$
iff $t\notin\operatorname{int}(\mathrm{NFP}(A,B))$; and $A$ and $B+t$ touch iff $t\in\partial\,\mathrm{NFP}(A,B)$.
\end{theorem}
\begin{proof}
$(B+t)\cap A\neq\varnothing$ iff $\exists\,a\in A,b\in B$ with $a=b+t$, i.e.\ $t=a-b\in A\oplus(-B)=\mathrm{NFP}(A,B)$.
Restricting to interiors, overlap of open interiors corresponds exactly to $t\in\operatorname{int}(\mathrm{NFP})$;
boundary contact (touching, no overlap) to $t\in\partial\,\mathrm{NFP}$. The NFP is computed by Clipper2's
polygon Minkowski sum (orbiting $-B$ around $A$).
\end{proof}
\begin{definition}[Inner-fit polygon / region]
$\mathrm{IFP}(S,B)=\{t:B+t\subseteq S\}$, the erosion of the sheet $S$ by $B$. For sheet \emph{holes} (defects)
the admissible region is $\mathrm{IFP}(S,B)\setminus\bigcup_h \mathrm{NFP}(H_h,B)$ and, for a part placed
\emph{inside} a host part's hole, $\mathrm{IFP}(\text{hole},B)$.
\end{definition}
\begin{definition}[Bottom-left-fill (BLF)]
Place parts in a fixed order; each part is positioned at the lexicographically smallest (bottom-most, then
left-most) $t$ in $\mathrm{IFP}(S,B)\setminus\bigcup_{k\,\text{placed}}\mathrm{NFP}(P_k,B)$.
\end{definition}
\begin{proposition}[BLF validity]
Every BLF placement satisfies $B+t\subseteq S$ (from IFP) and is interior-disjoint from all placed parts
(from removing each $\mathrm{NFP}$, Thm~\ref{thm:nfp}). Hence the layout is overlap-free by construction. (BLF
is a heuristic for density; optimality is not claimed.)
\end{proposition}
\noindent\emph{Hole-aware contact nesting (CNH).} The candidate translation set is sampled at NFP
\emph{contact} configurations and at edge-alignment rotations (the contact-adaptive angle set); holes are
nested host-first so smaller parts can seat inside larger parts' inner-fit regions.

\section{3D packing: orientations and deepest-left-bottom-fill}
\begin{definition}[Axis-aligned orientation group]
The proper rotation group of the cube (chiral octahedral group $O$) has order $24$. The 24-orientation set is
$\{R\in SO(3): R\,\text{maps the standard cube to itself}\}$, enumerated as the permutation--sign matrices
with $\det=+1$.
\end{definition}
\begin{lemma}
$|O|=24$. \emph{Proof:} a cube has $6$ faces that can map to the bottom and $4$ rotations fixing that face,
$6\cdot4=24$; equivalently $O\cong S_4$ acting on the four body diagonals.
\end{lemma}
\begin{definition}[DLBF on a heightmap]
Maintain a heightmap $h:\text{grid}\to\R$. For each part (in non-increasing volume order) and each of the 24
orientations, the drop position is the grid cell minimizing, lexicographically, (resting height, $y$, $x$),
where resting height $=\max$ of $h$ over the part footprint; place it and raise $h$ by the part's local
thickness. Equivalent to deepest-bottom-left-fill in 3D.
\end{definition}
\begin{proposition}[Non-overlap + support]
A DLBF placement rests on the current heightmap, so no two placed parts interpenetrate (each occupies cells
above the prior surface), and the part is supported (its base touches $h$). Mesh-accurate variants replace the
footprint-max by a CoACD collision test (\S Decomposition, future batch).
\end{proposition}

\section{Guillotine separability (wire-saw cutting)}
\begin{definition}[Guillotine cut and partition]
A \emph{guillotine cut} of an axis-aligned rectangle $R$ is a single edge-to-edge axis-parallel line splitting
$R$ into two sub-rectangles. A packing of blocks in $R$ is \emph{guillotine-partitionable} if $R$ can be
recursively reduced to the individual blocks by guillotine cuts.
\end{definition}
\begin{theorem}[Wire-saw separability $\Leftrightarrow$ guillotine]\label{thm:guillotine}
A block layout is separable by full-span (diamond-wire) cuts iff it is guillotine-partitionable; the staged
three-stage guillotine (bed cuts $\to$ strips $\to$ cross cuts) produces such a partition by construction, so
its separable fraction $\varphi=1$.
\end{theorem}
\begin{proof}
A diamond wire makes exactly full-span straight cuts $=$ guillotine cuts, so a layout is wire-separable iff
each freeing cut is edge-to-edge on the current piece, i.e.\ a guillotine sequence exists. The staged scheme
emits, at each stage, parallel full-span cuts on the current sub-region (stage 1 along beds, stage 2 into
strips, stage 3 cross), which is a recursive guillotine partition; hence every block is freed by full-span
cuts and $\varphi=1$.
\end{proof}
\begin{definition}[BCSdbBV cost, Jalalian]
For a recovered set, the cutting-surface objective is $\mathrm{I}_{11}=\dfrac{A_{\mathrm{cut}}}{\sum_b V_b\,\mathrm{price}_b}$,
the total sawn surface area per unit recovered block value; minimizing $\mathrm{I}_{11}$ minimizes saw work
(time/energy/waste) per value. Frahan adds it as an optional score term ($\mathrm{CutSurfaceWeight}$),
off by default, so it never alters the baseline packing.
\end{definition}

\section{Masonry equilibrium (CRA limit analysis)}
\begin{definition}[Rigid-block assembly and contacts]
An assembly is a set of rigid blocks $\{B_1,\dots,B_m\}$ with planar \emph{contact interfaces}; each interface
is discretized into contact points $j$ with local frame $(\nu_j,\tau_j^1,\tau_j^2)$ ($\nu_j$ the inward unit
normal). The contact force at $j$ is $f_j=f_j^n\nu_j+\sum_a f_j^{a}\tau_j^a\in\R^3$, $f_j^n$ the normal
(compression) magnitude.
\end{definition}
\begin{definition}[Equilibrium operator]
Stacking all contact forces into $f$ and the per-block dead loads (weights) into $g$, force/moment balance of
every block is the linear system $A f = g$, where $A$ is the equilibrium (statics) matrix assembled from the
contact positions and frames.
\end{definition}
\begin{definition}[Admissible cone: no-tension + Coulomb friction]
A contact force is admissible iff $f_j^n\ge0$ (unilateral: contacts push, never pull) and
$\Norm{(f_j^1,f_j^2)}\le\mu\,f_j^n$ (Coulomb friction cone, coefficient $\mu$). The admissible set
$\mathcal K=\{f: f_j^n\ge0,\ \Norm{(f_j^1,f_j^2)}\le\mu f_j^n\ \forall j\}$ is a product of second-order
(ice-cream) cones, hence closed and convex.
\end{definition}
\begin{theorem}[Static (safe) theorem of limit analysis]\label{thm:cra}
The assembly is stable under load $g$ iff there exists $f\in\mathcal K$ with $Af=g$. Equivalently, stability
is feasibility of the convex (second-order cone) program $\textsf{find } f\in\mathcal K \text{ s.t. } Af=g$.
\end{theorem}
\begin{proof}
This is the lower-bound (static) theorem of plastic limit analysis specialized to rigid no-tension blocks
with Coulomb interfaces (Heyman; Kao 2022 CRA): any statically admissible self-stress field in equilibrium
with the loads and respecting the yield set ($\mathcal K$) certifies that the structure can carry the load
(safe). Conversely, if no admissible equilibrium force exists, by convex (Gale/Farkas) duality there is a
compatible collapse mechanism with positive external work, so the assembly moves --- unstable. Convexity of
$\mathcal K$ and linearity of $A$ make this an SOCP feasibility test (CRA solves it; a polyhedral cone
linearization reduces it to an LP).
\end{proof}
\begin{definition}[Collapse load factor $\lambda$]
Split $g=g_{\mathrm{dead}}+\lambda\,g_{\mathrm{live}}$. The collapse multiplier is
$\lambda^\star=\sup\{\lambda: \exists f\in\mathcal K,\ Af=g_{\mathrm{dead}}+\lambda g_{\mathrm{live}}\}$; the
design is safe iff $\lambda^\star\ge1$. ($\lambda^\star$ is the optimum of a convex program; Frahan's
``Lambda'' stability index lives in this family and is used to compare assemblies, e.g.\ ETH1100 vs
Cyclopean.)
\end{definition}
\begin{proposition}[COM-over-support gate]
For a single rigid block resting on coplanar frictionless supports, the frictionless, single-interface case
of Thm~\ref{thm:cra} reduces to: the block does not topple iff the vertical projection of its centre of mass
lies in the convex hull of the support contact points.
\end{proposition}
\begin{proof}
Frictionless $\Rightarrow$ each $f_j=f_j^n\nu_j$ with $f_j^n\ge0$ and $\nu_j$ vertical. Force balance fixes
$\sum f_j^n=W$; moment balance about the horizontal axes requires the weighted average of support points
(weights $f_j^n\ge0$, summing to $W$) to equal the COM projection. A convex combination of the support points
realises a given point iff that point is in their convex hull. Hence feasibility $\Leftrightarrow$ COM
projection $\in\conv(\text{supports})$.
\end{proof}
\begin{definition}[Ashlar / coursed bond; Cyclopean carve]
A \emph{coursed ashlar} layout stacks blocks in horizontal courses with head joints of consecutive courses
offset by at least an overlap fraction $o$ (running bond: no two vertically adjacent head joints collinear),
which the layout engine enforces as a constraint maximizing interlock. \emph{Cyclopean Cannibalism} is the
overlap-then-carve operation: place an oversized boulder, then Boolean-subtract the already-placed neighbours
(carve-back) so faces mate exactly while preserving the rubble texture.
\end{definition}

\section{Decomposition}
\subsection{Approximate convex decomposition (CoACD)}
\begin{definition}[Concavity]
For a compact $S\subset\R^d$ with convex hull $\conv(S)$, the concavity is
$\eta(S)=\sup_{x\in\conv(S)}\operatorname{dist}(x,S)$ (the one-sided Hausdorff distance from $\conv(S)$ to
$S$). $\eta(S)=0$ iff $S$ is convex.
\end{definition}
\begin{definition}[$\varepsilon$-approximate convex decomposition]
A partition $\{S_k\}$ of $S$ (union $=S$, interiors disjoint) is an $\varepsilon$-ACD if $\eta(S_k)\le\varepsilon$
for all $k$. The decision ``$\exists$ $\varepsilon$-ACD with $\le K$ parts'' is NP-hard in 3D; CoACD is a
heuristic: a Monte-Carlo tree search over straight cut planes, scored by a \emph{collision-aware} concavity
proxy (it cuts the mesh and measures the post-cut concavity), recursing until every part has
$\eta\le\varepsilon$.
\end{definition}
\begin{proposition}[Termination + soundness]
Each CoACD cut strictly partitions a part with $\eta>\varepsilon$; since a finite triangle mesh admits only
finitely many distinct cut-induced sub-meshes and $\eta$ is non-increasing under cutting, the recursion
terminates with all parts $\eta\le\varepsilon$. The parts tile $S$ (union/disjoint-interior) by construction.
\end{proposition}

\subsection{Convex polygon partition}
\begin{theorem}[Hertel--Mehlhorn]\label{thm:hm}
Triangulate a simple polygon $P$ ($n$ vertices), then delete every \emph{inessential} diagonal (one whose
removal keeps both incident pieces convex). The result is a partition of $P$ into convex pieces whose count is
at most $4$ times the minimum, computed in $O(n)$ time after the triangulation.
\end{theorem}
\begin{proof}[Proof sketch]
A diagonal is essential only if it resolves a reflex vertex; each reflex vertex needs $\ge1$ resolving edge,
giving a lower bound, while Hertel--Mehlhorn uses $\le2$ essential diagonals per reflex vertex, yielding the
$4$-approximation (a reflex vertex is shared by $\le2$ kept diagonals). Deleting inessential diagonals never
creates a reflex angle, so all pieces stay convex.
\end{proof}
\noindent The exact minimum convex partition without Steiner points is solvable in $O(n^3)$ by dynamic
programming (Chazelle--Dobkin).

\subsection{Convex polyhedron half-space clipping}
\begin{lemma}[3D clip]\label{lem:clip3d}
For a convex polyhedron $Q\subset\R^3$ and half-space $H$, $Q\cap H$ is a convex polyhedron with
$Q\cap H\subseteq Q$; it is computed by clipping each face polygon against $H$ (planar Sutherland--Hodgman)
and capping the new cut face. This is the 3D analogue of Lemma~\ref{lem:sh} and the managed
half-space-clipping decomposer.
\end{lemma}
\begin{proof}
Intersection of convex sets is convex; $H\supseteq Q\cap H$ gives the subset. The boundary of $Q\cap H$
consists of clipped original faces plus one new face on $\partial H$ (the section), each a convex polygon.
\end{proof}

\subsection{Delaunay and $\alpha$-shapes}
\begin{definition}[Delaunay complex]
$\mathrm{DT}(P)$ is the simplicial complex on $P\subset\R^d$ in which a simplex belongs iff it has an
\emph{empty} circumsphere (no point of $P$ strictly inside). For points in general position it is the unique
triangulation maximizing the minimum angle (2D) / dual to the Voronoi diagram.
\end{definition}
\begin{definition}[$\alpha$-complex]
For $\alpha\ge0$, the $\alpha$-complex $\mathcal C_\alpha\subseteq\mathrm{DT}(P)$ keeps each simplex whose
smallest empty circumsphere has radius $\le\alpha$ (with the standard ``$\alpha$-exposed'' condition). The
$\alpha$-shape is $|\mathcal C_\alpha|$.
\end{definition}
\begin{theorem}[$\alpha$ interpolates points $\to$ hull]
$\mathcal C_0=P$ (vertices only) and for $\alpha\ge\alpha_{\max}$, $|\mathcal C_\alpha|=\conv(P)$ with
$\mathcal C_\alpha=\mathrm{DT}(P)$; the family $(\mathcal C_\alpha)_{\alpha\ge0}$ is a nested filtration that
grows from the point cloud to the convex hull. Constrained Delaunay tetrahedralization adds required facets to
$\mathrm{DT}$ while keeping the empty-circumsphere property on the unconstrained simplices.
\end{theorem}

\subsection{Density-watershed partition}
\begin{definition}[Watershed blocks (BlockCutOpt I5)]
Given a scalar fracture-density field $f$ on a grid, define $g=-f$ and assign each cell to the catchment
basin of the local maximum of $g$ (= local minimum of $f$) reached by steepest ascent; the basins (separated
by watershed ridges along high-density fractures) partition the volume into candidate intact blocks.
\end{definition}
\begin{proposition}
The catchment-basin map is a partition: every non-ridge cell flows to exactly one maximum, and basins are
interior-disjoint and cover the domain up to the measure-zero ridge set. Low fracture density $\Rightarrow$
large intact basins $\Rightarrow$ larger recoverable blocks.
\end{proposition}

\section{Edge-matching and assembly}
\subsection{Closed-form absolute orientation (Horn 1987)}
\begin{definition}[Absolute orientation]
Given corresponded point sets $\{p_i\},\{q_i\}_{i=1}^n\subset\R^3$, find $s>0$, $R\in SO(3)$, $t\in\R^3$
minimizing $E(s,R,t)=\sum_i\Norm{q_i-(sRp_i+t)}^2$.
\end{definition}
\begin{theorem}[Horn quaternion solution]\label{thm:horn}
Let $\bar p,\bar q$ be the centroids, $p_i'=p_i-\bar p$, $q_i'=q_i-\bar q$, and
$M=\sum_i p_i' q_i'^\top$. Then the optimal rotation is $R=R(\hat e)$ where $\hat e$ is the unit eigenvector of
the largest eigenvalue of the symmetric $4\times4$ matrix $N(M)$ (built from the entries of $M$); the optimal
scale is $s=\big(\sum_i\Norm{q_i'}^2/\sum_i\Norm{p_i'}^2\big)^{1/2}$ (symmetric form) and $t=\bar q-sR\bar p$.
\end{theorem}
\begin{proof}[Proof sketch]
Minimizing over $t$ for fixed $s,R$ gives $t=\bar q-sR\bar p$ (centroid alignment), reducing $E$ to the
centred residual. The rotation part maximizes $\sum_i q_i'^\top R p_i'=\operatorname{tr}(R M^\top)$. Writing
$R$ via a unit quaternion $\dot e$, this trace equals $\dot e^\top N(M)\dot e$ with $N(M)$ the standard
symmetric $4\times4$ form; maximizing a quadratic form over the unit sphere is attained at the top eigenvector
(Rayleigh). Scale follows by setting $\partial E/\partial s=0$. (Equivalent SVD form: $R=V\,\mathrm{diag}(1,1,\det VU^\top)U^\top$ from $M=U\Sigma V^\top$.)
\end{proof}

\subsection{Iterative closest point (ICP)}
\begin{definition}
Given a source $\{p_i\}$ and target set $Y$, ICP alternates: (C) $c(i)=\argmin_{y\in Y}\Norm{T p_i-y}$
(nearest correspondence under the current transform $T$); (T) $T\leftarrow\argmin_{T}\sum_i\Norm{Tp_i-c(i)}^2$
(Horn, Thm~\ref{thm:horn}). \emph{Constrained} ICP adds non-penetration/contact constraints to step (T).
\end{definition}
\begin{theorem}[Monotone convergence]
The registration error $E_t=\sum_i\Norm{T_t p_i-c_t(i)}^2$ is non-increasing in $t$ and converges; ICP
reaches a fixed point (local minimum) in finitely many correspondence sets.
\end{theorem}
\begin{proof}
Step (C) minimizes $E$ over correspondences with $T$ fixed (each $c(i)$ is the nearest point), so it cannot
increase $E$; step (T) minimizes over $T$ with correspondences fixed, so it cannot increase $E$. Thus $E_t$ is
monotone non-increasing and bounded below by $0$, hence convergent; since correspondences come from a finite
set and $E$ strictly decreases off a fixed point, ICP terminates at a local minimum.
\end{proof}

\subsection{Global assembly: matching, beam, spanning tree}
\begin{definition}[Bipartite assignment]
Given a cost matrix $C\in\R^{n\times n}$, the assignment problem is
$\min_{\pi\in S_n}\sum_i C_{i,\pi(i)}$, solved exactly by the Hungarian algorithm in $O(n^3)$ (it returns a
minimum-cost perfect matching by maintaining a dual feasible potential and augmenting along tight edges).
\end{definition}
\begin{definition}[Pair-graph assembly]
Build a weighted graph $G$ on fragments with edge weight the pairwise match score (boundary-rail edge
affinity). \emph{Agglomerative} assembly takes a maximum-weight spanning tree of $G$ and composes the
relative poses along it; \emph{beam search} keeps the top-$B$ partial assemblies by cumulative score, expanding
each by one compatible fragment per step. Beam with $B=1$ is the frame-anchored greedy; $B\to\infty$ is exact.
\end{definition}
\begin{proposition}[Spanning-tree consistency]
A spanning tree induces a unique relative pose between any two fragments (compose along the tree path); cycles
are avoided, so the assembly is pose-consistent by construction (no contradictory loop closures), at the cost
of not averaging redundant matches.
\end{proposition}

\section{Surface reconstruction and flattening}
\subsection{Poisson surface reconstruction (Kazhdan)}
\begin{definition}[Indicator and oriented-normal field]
Let $M=\partial\Omega$ be the surface of a solid $\Omega$ with indicator $\chi_\Omega$. The oriented-point
samples define a vector field $V$ approximating the surface normal field as a measure: $\nabla\chi_\Omega=\vec
n\,\delta_M$ in the distributional sense.
\end{definition}
\begin{theorem}[Poisson formulation]\label{thm:poisson}
The indicator $\chi$ best matching $V$ in the least-squares sense solves the Poisson equation
$\Delta\chi=\nabla\cdot V$, and $M$ is recovered as a level set $\{\chi=\bar\chi\}$ (isovalue = mean of $\chi$
at the samples).
\end{theorem}
\begin{proof}[Sketch]
Minimizing $\int\Norm{\nabla\chi-V}^2$ over $\chi$ has Euler--Lagrange equation $\Delta\chi=\nabla\cdot V$ (set
the variational derivative to zero, integrate by parts). Since $\nabla\chi_\Omega=\vec n\,\delta_M$ and $V$
approximates that field, $\chi\approx\chi_\Omega$, whose level set at the surface is $M$. (Discretized on an
adaptive octree; the right-hand side is the divergence of the smoothed normal field.)
\end{proof}

\subsection{Centroidal Voronoi tessellation (Lloyd)}
\begin{definition}[CVT energy]
For generators $z=(z_1,\dots,z_k)$ with Voronoi cells $V_i(z)$ and density $\rho$, the energy is
$F(z)=\sum_i\int_{V_i(z)}\rho(x)\Norm{x-z_i}^2\,dx$. A tessellation is \emph{centroidal} if each $z_i$ equals
its mass centroid $c_i=\big(\int_{V_i}\rho x\big)/\big(\int_{V_i}\rho\big)$.
\end{definition}
\begin{theorem}[Lloyd descent]
Lloyd's iteration (recompute $V_i(z)$; set $z_i\leftarrow c_i$) does not increase $F$ and its fixed points are
exactly the centroidal Voronoi tessellations.
\end{theorem}
\begin{proof}
With cells fixed, $\sum_i\int_{V_i}\rho\Norm{x-z_i}^2$ is minimized over $z_i$ at the centroid $c_i$ (the
mean minimizes squared deviation), so the move step lowers $F$; re-Voronoi with generators fixed assigns each
$x$ to its nearest $z_i$, lowering each integrand pointwise, so it also lowers $F$. $F\ge0$ bounded below
$\Rightarrow$ convergence; a fixed point has $z_i=c_i\,\forall i$, i.e.\ a CVT. (Used for Trencadis seeding and
CVD remeshing.)
\end{proof}

\subsection{Barycentric mapping and conformal flattening}
\begin{definition}[Barycentric map]
For a triangle $T=(a,b,c)$, the map $\varphi_T(\lambda)=\lambda_1 a+\lambda_2 b+\lambda_3 c$ on the standard
simplex $\{\lambda\ge0,\sum\lambda=1\}$ is the unique affine bijection from the reference triangle to $T$.
\end{definition}
\begin{proposition}[2D$\to$3D transfer]
A point with barycentric coordinates $\lambda$ in a parameter triangle maps to $\varphi_{T_{3D}}(\lambda)$ in
the corresponding 3D triangle. Barycentric coordinates are affine invariants, so the transfer is independent
of the embedding and exact on triangles.
\end{proposition}
\begin{definition}[BFF flattening]
Boundary First Flattening computes a discrete \emph{conformal} map of a disk-topology surface to the plane by
prescribing either boundary lengths or boundary geodesic curvatures and solving the resulting Cherrier/Yamabe
boundary-value problem (a sparse linear / Poisson solve for the log conformal factor $u$, then a boundary
integration). Conformality $\Rightarrow$ angles preserved, area distortion encoded by $e^{2u}$ (the
chart-scale, recovered for the surface packers).
\end{definition}

\section{Discontinuities, joint sets, and block theory}
\subsection{Baecher stochastic DFN}
\begin{definition}
A Baecher discrete fracture network draws fracture centres from a homogeneous Poisson point process of
intensity $\lambda$ in the rock volume; each fracture is a disk of radius $r\sim F_R$ (size law) with unit
normal $\nu\sim$ Fisher on the sphere, density $f(\nu)\propto\exp(\kappa\,\mu^\top\nu)$ ($\mu$ mean pole,
$\kappa$ concentration). Centres, sizes, orientations are independent.
\end{definition}
\subsection{Watson axial mean-shift for joint sets}
\begin{definition}
Fracture poles are \emph{axial} ($\nu\equiv-\nu$). The Watson distribution on the projective sphere has
density $f(x)\propto\exp(\kappa(\mu^\top x)^2)$ (bipolar for $\kappa>0$). Joint sets are the modes of the pole
density, found by \emph{mean shift} on $\R P^2$: iterate $x\leftarrow$ normalize$\big(\sum_i w_i(x)\,
\mathrm{sign}(x^\top\nu_i)\,\nu_i\big)$ with kernel weights $w_i(x)=g(\,1-(x^\top\nu_i)^2\,)$.
\end{definition}
\begin{proposition}
Each mean-shift step ascends the kernel pole-density estimate; fixed points are the density modes, partitioning
the poles into joint sets (basins of attraction). Antipodal symmetry is handled by the squared inner product,
so $\nu$ and $-\nu$ vote together.
\end{proposition}
\subsection{Block theory: removability (Goodman--Shi)}
\begin{definition}[Joint, excavation, block pyramids]
For a convex block cut by joint planes through a common apex, the \emph{joint pyramid} $JP$ is the cone
$\{d:\ \text{$d$ on the block side of every joint plane}\}$; the \emph{excavation pyramid} $EP$ is the cone of
directions into the free (excavated) space; the \emph{block pyramid} is $BP=JP\cap EP$.
\end{definition}
\begin{theorem}[Finiteness and removability]\label{thm:blocktheory}
A block is \emph{finite} iff $BP=JP\cap EP=\varnothing$. A finite block is \emph{removable} iff $JP\neq
\varnothing$; if instead $JP=\varnothing$ it is \emph{tapered} (cannot be removed without cutting). Hence
\[\text{removable}\iff JP\neq\varnothing\ \wedge\ JP\cap EP=\varnothing.\]
A removable block is a \emph{key block} iff the active resultant force lies in $JP$ (so it slides/falls under
gravity/water/seismic load).
\end{theorem}
\begin{proof}[Sketch]
Shi's theorem: the block translated to infinity along a direction $d$ stays in the rock iff $d\in JP$ and in
the free space iff $d\in EP$; an unbounded motion (removal without cutting) requires a common direction, i.e.\
$d\in BP$. Thus $BP=\varnothing\Rightarrow$ no escape to infinity $\Rightarrow$ finite, and a finite block can
be moved (removable) iff its shape cone $JP$ is non-empty (some sliding/lifting direction exists once the
excavation face is open), while $JP=\varnothing$ forces a tapered, non-removable wedge. Convex-cone
intersection tests make this a linear-feasibility check per candidate block.
\end{proof}
\subsection{Block size from joint spacing (Palmstrom)}
\begin{definition}
With volumetric joint count $J_v=\sum_i 1/s_i$ ($s_i$ mean spacing of joint set $i$), the typical block volume
is $V_b=\beta\,J_v^{-3}$ ($\beta$ a block-shape factor). $\mathrm{CrackGraph}\to\mathrm{BlockGraph}$ partitions
the fractured solid into the cells of the joint arrangement, whose volumes realize this statistic.
\end{definition}

\section{GPR signal processing (pre-kriging)}
\subsection{f--k (Stolt) migration}
\begin{definition}[Exploding-reflector section]
A zero-offset radargram $P(x,t)$ is modeled as an exploding-reflector field with effective velocity
$c=v/2$ (half the EM velocity), so reflectors radiate at $t=0$.
\end{definition}
\begin{theorem}[Stolt mapping, constant $v$]\label{thm:stolt}
Let $\hat P(k_x,\omega)$ be the 2D Fourier transform of $P$. The migrated image $M(x,z)$ has transform
$\hat M(k_x,k_z)=\dfrac{c\,k_z}{\sqrt{k_x^2+k_z^2}}\,\hat P\!\big(k_x,\ \omega=c\sqrt{k_x^2+k_z^2}\big)$, i.e.\
a remap from $\omega$ to $k_z=\sqrt{(\omega/c)^2-k_x^2}$ with the indicated amplitude Jacobian, followed by an
inverse 2D FFT.
\end{theorem}
\begin{proof}[Sketch]
The scalar wave equation gives the dispersion relation $\omega^2=c^2(k_x^2+k_z^2)$ for upgoing waves; imaging
at $t=0$ evaluates the field at depth $z$. Changing the integration variable from $\omega$ to $k_z$ along
$\omega=c\sqrt{k_x^2+k_z^2}$ introduces the Jacobian $d\omega/dk_z=c\,k_z/\sqrt{k_x^2+k_z^2}$; the inverse FFT
then yields $M(x,z)$. This is exact for constant $c$ (Stolt 1978).
\end{proof}

\subsection{Analytic signal and instantaneous energy (Taner)}
\begin{definition}[Hilbert transform, analytic signal]
The Hilbert transform $\mathcal H$ has frequency multiplier $-i\,\mathrm{sgn}(\omega)$. The analytic signal of a
trace $s$ is $s_a=s+i\,\mathcal H s$; the instantaneous amplitude is $A(t)=|s_a(t)|$ and the instantaneous
energy $A(t)^2$ (Taner 1979).
\end{definition}
\begin{proposition}
$A(t)$ is the envelope of $s$: $A\ge|s|$ with equality at the local extrema, and $A$ is phase-independent
(invariant under a constant phase rotation $s_a\mapsto e^{i\phi}s_a$). Hence energy peaks localize reflectors
robustly to wavelet phase.
\end{proposition}

\subsection{Picking, continuity gate, depth, uncertainty}
\begin{definition}[Energy-quantile pick + USGS continuity]
A sample is a fracture candidate iff its instantaneous energy exceeds the $Q$-quantile of the trace energy.
A candidate reflector is \emph{kept} iff it is laterally continuous over at least $N$ consecutive traces
(USGS $\ge N$-trace gate; $N$ stone-calibrated, e.g.\ 27 marble vs 40 granite).
\end{definition}
\begin{proposition}[Depth and its uncertainty]
Two-way time $t$ maps to depth $d=v\,t/2$. With $v=c_0/\sqrt{\varepsilon_r}$, error propagation gives the
dominant deep-fracture deviation
\[\frac{\sigma_d}{d}=\frac{\sigma_v}{v}=\tfrac12\,\frac{\sigma_{\varepsilon_r}}{\varepsilon_r},\]
plus a $\lambda/4$ vertical-resolution floor and a time-zero term; these combine (root-sum-square) into the
per-pick $\sigma$ consumed by the kriging confidence metric (\S1, FractureUncertainty).
\end{proposition}

\section{Quarry yield and cutting optimization}
\subsection{Catalogue guillotine tiling (cost/volume)}
\begin{definition}[Priced catalogue, layer]
A catalogue $\mathcal B$ assigns each block size $b$ a footprint and a unit price/value. A bed-bounded layer is
a rectangle $R$; a guillotine tiling places non-overlapping catalogue blocks separable by edge-to-edge cuts.
The objective with volume weight $W$ is $\max\ \sum_b \mathrm{vol}_b(\mathrm{price}_b+W)\,-\,\kappa\,A_{\mathrm{cut}}$.
\end{definition}
\begin{theorem}[DP optimality among guillotine tilings]\label{thm:guillodp}
Let $V(R)$ be the optimum objective over guillotine tilings of $R$. Then
\[V(R)=\max\Big(\max_{b\subseteq R}\mathrm{val}(b)+V(R\setminus b),\ \max_{\text{cut }\ell}V(R_1)+V(R_2)\Big),\]
over single-block placements and all horizontal/vertical cut positions $\ell$ splitting $R=R_1\cup R_2$. The
recursion (memoized over the $O(n^2)$ distinct sub-rectangles induced by the catalogue/grid coordinates)
computes the exact guillotine optimum.
\end{theorem}
\begin{proof}
A guillotine tiling of $R$ is, by definition, either a single block or a first edge-to-edge cut into $R_1,R_2$
each guillotine-tiled; the value is additive over the partition; taking the max over the (finite) cut
positions and placements gives Bellman optimality. Restricting cuts to catalogue/grid coordinates loses no
optimum because shifting a cut to the nearest such coordinate does not reduce the achievable value.
\end{proof}
\noindent\emph{RecoveryCascade} applies $V$ at decreasing scales (coarse blocks first, then fill); at a
single scale it reduces to Thm~\ref{thm:guillodp}, and the multi-scale recovery is monotone (finer passes only
add blocks in the residual region).

\subsection{Pareto omni-solve (BlockCutOpt)}
\begin{definition}[Dominance, Pareto front]
For objective vector $\phi=(\text{recovery},\text{revenue},-\text{risk},-A_{\mathrm{cut}})$ (all maximized),
$x$ \emph{dominates} $y$ iff $\phi(x)\ge\phi(y)$ componentwise with strict inequality somewhere. The Pareto
front is the set of non-dominated solutions. BlockCutOpt omni-solve returns this front; the single-objective
solvers (brute force, density-watershed I5, Fisher-robust) are scalarizations / special cases.
\end{definition}
\begin{proposition}
The recovery-optimal and revenue-optimal plans both lie on the Pareto front but generally differ (\S1
finding); exposing the front lets the operator pick the trade-off rather than a fixed scalarization.
\end{proposition}

\subsection{Cut-and-fill volume (TIN prisms)}
\begin{proposition}[Prism differencing]
For two TINs (top, bottom) over a common triangulation, the volume between them is
$V=\sum_{T}\dfrac{A_T}{3}\big(\delta_{1}+\delta_2+\delta_3\big)$, where $A_T$ is the planimetric area of
triangle $T$ and $\delta_i$ the top-minus-bottom height at its vertices (truncated-prism / prismoidal rule);
signed $\delta$ separates cut ($\delta<0$) from fill ($\delta>0$).
\end{proposition}
\begin{proof}
Over a triangle the height difference $\delta(x,y)$ is affine (linear TINs), so its mean over $T$ equals the
vertex average $(\delta_1+\delta_2+\delta_3)/3$; the prism volume is mean height times base area $A_T$. Sum
over triangles.
\end{proof}

\section{Robotic fabrication: toolpaths and frames}
\subsection{Tool-axis frame along a cut path}
\begin{definition}[Path frame]
For a cut path $\gamma(s)$ (arc length) with unit tangent $T=\gamma'$, a tool frame is an orthonormal
$(T,U,V)$; the tool axis is one of $U,V$ (or the cut-plane normal). A \emph{rotation-minimizing} (Bishop)
frame parallel-transports $U,V$: $U'=-(U\cdot T')\,T$, $V'=-(V\cdot T')\,T$.
\end{definition}
\begin{proposition}[RMF minimizes spin]
The Bishop frame has zero angular velocity component along $T$ ($\Omega\cdot T=0$), hence among all adapted
frames it minimizes total tool roll $\int|\Omega\cdot T|\,ds=0$; it exists and is unique given an initial
$(U_0,V_0)\perp T(0)$. (Avoids the Frenet frame's flips at inflections, giving smooth robot wrist motion.)
\end{proposition}
\begin{proof}
$U'=-(U\cdot T')T$ keeps $U\perp T$ ($\frac{d}{ds}(U\cdot T)=U'\cdot T+U\cdot T'=-(U\cdot T')+U\cdot T'=0$) and
$\Norm U=1$ ($U\cdot U'=0$). Its angular velocity is $\Omega=T\times T'$, which has no $T$-component, so the
roll rate $\Omega\cdot T=0$. Uniqueness/existence is the linear ODE with given initial frame.
\end{proof}
\subsection{Arc discretization (Arc Step) and motion mapping}
\begin{proposition}[Chord-error bound]
An arc of radius $r$ and sweep $\theta$ split into $n$ equal chords (each subtending $\alpha=\theta/n$) has
maximum deviation (sagitta) $h=r\big(1-\cos(\alpha/2)\big)$. To guarantee $h\le\varepsilon$ take
$\alpha\le2\arccos(1-\varepsilon/r)$, i.e.\ $n=\big\lceil \theta/\big(2\arccos(1-\varepsilon/r)\big)\big\rceil$.
\end{proposition}
\begin{proof}
The chord's farthest point from the arc is its midpoint; the perpendicular distance from the arc to the chord
midpoint is $r-r\cos(\alpha/2)$ (radius minus the apothem). Bounding it by $\varepsilon$ and solving for
$\alpha$ gives the stated $n$.
\end{proof}
\begin{definition}[CutSegment $\to$ motion morphism]
Map each cut segment kind to a robot motion primitive: line$\to$\textsc{Lin}, arc$\to$\textsc{Circ} (via three
points), rapid$\to$\textsc{Ptp}; the tool pose at each via point is $(\gamma(s_i),\,\text{frame}(s_i))$. The map
is a structure-preserving morphism: the executed TCP path equals the geometric cut path up to the discretized
chord tolerance above (KUKA\,prc / Robots back-ends).
\end{definition}

\section{Rigid-body dynamics and settling}
\subsection{Newton--Euler equations of motion}
\begin{definition}
A rigid body with mass $m$, inertia $I$, position $x$, orientation $R$, linear/angular velocity $v,\omega$
evolves by $\dot x=v$, $\dot R=[\omega]_\times R$, and the Newton--Euler equations
$m\dot v=F$, $I\dot\omega+\omega\times I\omega=\tau$, with $F=mg+\sum_c f_c$, $\tau=\sum_c r_c\times f_c$
($f_c$ the contact forces, $r_c$ the contact arms).
\end{definition}
\subsection{Contact as a complementarity (Signorini) problem}
\begin{definition}[Signorini--Coulomb contact]
At each contact let $\phi$ be the signed gap and $\lambda$ the normal impulse. Non-penetration and
unilateral force are the complementarity condition $0\le\phi\ \perp\ \lambda\ge0$ (with the Coulomb friction
cone, \S6). Per step this is a (linear) complementarity problem; the sequential-impulse / projected
Gauss--Seidel solver iterates each contact to satisfy it.
\end{definition}
\begin{proposition}[Collision via convex decomposition]
For non-convex bodies decomposed into convex pieces (\S Decomposition, CoACD), interpenetration holds iff some
pair of pieces overlaps; each convex pair is tested by GJK (distance) / EPA (penetration). Thus narrow-phase
collision reduces to finitely many convex queries, each solvable by Minkowski-difference support mapping.
\end{proposition}
\subsection{Drop-settle and static rest}
\begin{theorem}[Rest = constrained energy minimum = contact equilibrium]\label{thm:settle}
A dropped block is at static rest iff its pose is a KKT point of
$\min_{q}\ U(q)=m g^\top c(q)$ subject to the non-penetration constraints $\phi_j(q)\ge0$, where $c(q)$ is the
centre of mass. The KKT stationarity $\nabla U=\sum_j\lambda_j\nabla\phi_j$, $\lambda_j\ge0$,
$\lambda_j\phi_j=0$ is exactly the contact-force equilibrium of \S6 (Thm~\ref{thm:cra}); hence
drop-settle (gradient descent of $U$ under the unilateral constraints, $z$-up) converges to a locally stable
configuration, and the resulting support reaction set certifies stability via the COM-over-support gate.
\end{theorem}
\begin{proof}[Sketch]
At rest, velocities and accelerations vanish, so Newton--Euler reduces to force/torque balance under gravity
and contact reactions $\lambda_j\nabla\phi_j$ ($\lambda_j\ge0$ since contacts push) with $\lambda_j\phi_j=0$
(no force at open gaps). These are precisely the stationarity + complementarity KKT conditions of the
gravitational-PE minimization under $\phi_j\ge0$. Gradient flow of $U$ projected onto the feasible set
decreases $U$ monotonically and is bounded below (the body rests on the support), so it converges to such a
KKT point. Concave-aware settling uses the convex-piece contacts for $\phi_j$.
\end{proof}

\section{Learned reassembly (Kintsugi / PotNet)}
\subsection{Normalization and pose composition}
\begin{definition}[Normalization]
Each fragment $f$ is mapped to a canonical frame by a similarity $T_{\mathrm{norm}}(f)\in\mathrm{Sim}(3)$
(translate to centroid, scale to unit extent); its inverse is $T_{\mathrm{unnorm}}(f)=T_{\mathrm{norm}}(f)^{-1}$.
The network $T_{\mathrm{net}}(f)\in SE(3)$ predicts the placement of the \emph{normalized} fragment into the
normalized assembly frame anchored at fragment $0$.
\end{definition}
\begin{theorem}[World-pose composition]\label{thm:kintsugi}
The world pose of fragment $f$ is
\[T_{\mathrm{world}}(f)=T_{\mathrm{unnorm}}(0)\,\cdot\,T_{\mathrm{net}}(f)\,\cdot\,T_{\mathrm{norm}}(f),\]
and this is the unique transform consistent with: normalize $f$, apply the predicted placement, then undo the
anchor (fragment $0$) normalization. It is invariant to the global pose of the raw input (the network sees
only normalized geometry).
\end{theorem}
\begin{proof}
For a point $p$ on fragment $f$: normalization gives $T_{\mathrm{norm}}(f)p$; the network places it in the
normalized assembly frame as $T_{\mathrm{net}}(f)T_{\mathrm{norm}}(f)p$; mapping the normalized assembly frame
back to world is the inverse of the anchor normalization, $T_{\mathrm{unnorm}}(0)$. Composing (associativity of
$SE(3)/\mathrm{Sim}(3)$ action) yields $T_{\mathrm{world}}(f)p$. If the raw input is pre-transformed by any
$G$, then $T_{\mathrm{norm}}$ absorbs $G$ (centroid/extent are equivariant), so $T_{\mathrm{net}}$ and the
final world relation are unchanged --- global-pose invariance. (PotNet regresses one such $T_{\mathrm{net}}$
per object; failure modes separate into pose-composition vs network-drift vs mesh-style by inspecting the
verifier-score distribution.)
\end{proof}
\subsection{Verifier-score gating}
\begin{definition}
A learned verifier $g:\text{assembly}\to[0,1]$ scores a candidate; it is accepted iff $g\ge\tau$
($\tau=0.5$). The accepted set is $\{f: g(f)\ge\tau\}$, a thresholding decision whose precision/recall trade
off monotonically as $\tau$ varies (higher $\tau$ $\Rightarrow$ higher precision, lower recall).
\end{definition}

\section{Surface packing on charts}
\subsection{Atlas and pack-then-map}
\begin{definition}[Face-corner UV chart]
A chart assigns each face corner a UV coordinate; the chart is a valid atlas iff the UV assignment is
consistent across every shared edge (the two incident faces agree up to the transition map). For a BFF flattening
(\S15) the chart is a single global UV map of the disk-topology surface.
\end{definition}
\begin{theorem}[Surface-packing transfer]\label{thm:surfpack}
Let parts be packed without overlap in the UV chart $\Omega$ (the 2D nesting of \S2). Mapping each placed part
back by the chart map $\Phi=\varphi\circ(\text{chart})^{-1}$ yields surface-embedded parts that are mutually
non-overlapping on the surface, and whose physical area equals the UV area times the local chart-scale
$e^{2u}$ (the BFF conformal factor); pre-scaling each UV cell by $e^{-u}$ makes the on-surface tile sizes
match the target.
\end{theorem}
\begin{proof}
$\Phi$ is a bijection on the chart (atlas validity), so disjoint UV placements map to disjoint surface
placements (\S2 non-overlap is preserved by a bijection). For a conformal chart the first fundamental form is
$e^{2u}\,\mathrm{Id}$, so an area element $dA_{\mathrm{UV}}$ maps to $e^{2u}dA_{\mathrm{UV}}$ on the surface;
hence multiplying UV sizes by $e^{-u}$ cancels the distortion (barycentric transfer, \S15, is exact per
triangle).
\end{proof}
\subsection{Orientation field (GVF)}
\begin{definition}[Gradient vector flow field]
Given an initial tile-orientation field $v_0$ (e.g.\ from surface curvature / edges), the smoothed field $v$
minimizes the Xu--Prince energy $E(v)=\int_\Omega \mu\,\Norm{\nabla v}^2+\Norm{v_0}^2\,\Norm{v-v_0}^2\,dA$.
\end{definition}
\begin{proposition}[Well-posed orientation]
$E$ is a strictly convex quadratic functional in $v$ (for $\mu>0$), so it has a unique minimizer solving the
linear diffusion--reaction Euler--Lagrange equation $\mu\,\Delta v-\Norm{v_0}^2(v-v_0)=0$; the result is a
smooth orientation field that follows $v_0$ where it is strong and diffuses smoothly elsewhere, used to align
Trencadis mosaic tiles. CVD-Lloyd seeding (\S15) places the tile centres.
\end{proposition}
\noindent\emph{Coverage.} The remaining \texttt{[Algorithm]} attributes that are genuine mathematical methods
(not format I/O or dispatch) are derived in \S20--\S27 below; see \texttt{AUDIT\_coverage.md} for the
title-by-title mapping of all 249 distinct algorithms.

\section{Spectral and PCA geometry}
\begin{definition}[PCA]
For points $\{p_i\}$ with centroid $\bar p$, the covariance is $C=\tfrac1n\sum_i(p_i-\bar p)(p_i-\bar p)^\top$,
with eigenpairs $(\lambda_1\ge\lambda_2\ge\lambda_3,\,u_1,u_2,u_3)$. The $u_k$ are the principal axes.
\end{definition}
\begin{proposition}[PCA-OBB and normal estimation]\label{prop:pca}
(i) The PCA oriented bounding box is the axis-aligned box of $\{p_i\}$ in the eigenbasis $U=[u_1u_2u_3]$; it is
a heuristic (not minimum-volume, which needs rotating calipers) but tight when one axis dominates.
(ii) The surface normal at a point is $u_3$ (least-eigenvalue direction) of the local-neighbourhood covariance
(cf.\ \S5). (iii) Hoppe normal orientation propagates a consistent sign along the Euclidean minimum spanning
tree, maximizing $\sum_{(i,j)\in MST}\mathrm{sign}(u_3^{(i)}\!\cdot u_3^{(j)})$.
\end{proposition}
\begin{proof}
$C\succeq0$ symmetric, so the spectral theorem gives the orthonormal $u_k$; the box in coordinates $U^\top p$
is the OBB. The least-variance direction minimizes $u^\top C u$ over $\Norm u=1$ (Rayleigh) and is normal to
the best-fit plane (\S5). MST propagation is the greedy consistent-sign assignment minimizing flips on a tree
(unique path between any two nodes $\Rightarrow$ no conflict).
\end{proof}

\section{Convex hull and inscribed polygons}
\begin{proposition}[QuickHull]
QuickHull computes $\conv(S)$ by recursively assigning points to the outside of the current hull edge and
promoting the farthest; it is exact, expected $O(n\log n)$, worst case $O(n^2)$ --- an alternative to the
$O(n\log n)$ monotone chain of \S4 (same output).
\end{proposition}
\begin{definition}[Largest inscribed convex polygon / potato-peeling]
For a simple polygon $P$, the \emph{convex skull} is the maximum-area convex polygon $Q\subseteq P$. The
restricted version fixes the vertex count (e.g.\ largest inscribed quadrilateral).
\end{definition}
\begin{theorem}[Exact convex blank vs.\ greedy trim]\label{thm:potato}
The convex skull is the area-optimal convex blank cut from a slab. The greedy half-plane trim of \S7
(Thm~\ref{thm:trim}) returns \emph{some} convex $Q\subseteq P$ (feasible, few cuts) but generally
sub-optimal; the convex skull is its area upper bound. Chang--Yap solves the convex skull exactly in $O(n^7)$;
a $(1-\varepsilon)$ near-linear approximation exists (Cabello et al.). Thus
$\area(\text{greedy }Q)\le\area(\text{convex skull})\le\area(P)$.
\end{theorem}
\begin{proof}
Both are convex subsets of $P$, so both have area $\le\area(P)$. The convex skull maximizes area over all
convex subsets by definition, dominating the greedy feasible solution. Exactness/complexity are the
Chang--Yap dynamic program over the polygon's visibility structure.
\end{proof}

\section{Voronoi family, power diagrams, offsets}
\begin{definition}[Power (Laguerre) diagram]
With sites $z_i$ and weights $w_i$, the power cell of $i$ is $\{x:\Norm{x-z_i}^2-w_i\le\Norm{x-z_j}^2-w_j\ \forall j\}$.
\end{definition}
\begin{proposition}[Convex polytope cells]\label{prop:power}
Each power cell is a (possibly empty) convex polytope: the defining inequalities are linear in $x$
(the quadratic $\Norm x^2$ cancels), $2(z_j-z_i)^\top x\le(\Norm{z_j}^2-w_j)-(\Norm{z_i}^2-w_i)$, an
intersection of half-spaces. The ordinary Voronoi diagram is the case $w_i\equiv0$; the restricted Voronoi
diagram intersects each cell with the surface mesh.
\end{proposition}
\begin{definition}[Straight skeleton and kerf offset]
The straight skeleton of a polygon is the trace of vertices under the \emph{mitered} inward offset (each edge
moves inward at unit speed); it partitions the polygon by straight segments (unlike the medial axis, no
parabolic arcs). A kerf offset of a cut curve is the Minkowski sum/erosion with a disk of radius
$\mathrm{kerf}/2$.
\end{definition}

\section{Mesh simplification, segmentation, geodesics}
\begin{definition}[Quadric error (Garland--Heckbert)]
Assign each vertex the quadric $Q=\sum_{p\in\text{planes}(v)} K_p$, $K_p=\begin{psmallmatrix}n\\d\end{psmallmatrix}
\begin{psmallmatrix}n\\d\end{psmallmatrix}^\top$ for plane $n^\top x+d=0$. The cost of contracting an edge to
$\bar v$ is $\bar v^\top(Q_1+Q_2)\bar v$ (homogeneous $\bar v$).
\end{definition}
\begin{theorem}[QEM = sum of squared plane distances]\label{thm:qem}
$v^\top K_p v=(n^\top \tilde v+d)^2=\operatorname{dist}(v,p)^2$, so $v^\top Q v$ is the sum of squared
distances from $v$ to its incident planes; the optimal contraction target solves the linear system
$\nabla(v^\top Q v)=0$, i.e.\ $\bar Q\,\bar v=-b$ (the upper $3\times3$ block $\bar Q$ with $b$ the linear
part), and greedy least-cost edge collapse gives the simplified mesh. Vertex-clustering decimation instead
snaps each grid cell to one representative, with Hausdorff error bounded by the cell diagonal.
\end{theorem}
\begin{proof}
Expand $v^\top K_p v$ with $\tilde v=(v,1)$: it equals $(n^\top v+d)^2$, the squared point-plane distance.
Summation gives $v^\top Q v$. $v^\top Q v$ is a convex quadratic; its minimizer is the stationary point
$\bar Q\bar v=-b$ when $\bar Q$ is invertible (else fall back to the edge midpoint).
\end{proof}
\begin{proposition}[Variational Shape Approximation]
VSA partitions faces into $k$ regions each with a plane proxy, alternating (proxy = area-weighted average
normal/point) and (region grow by $L^{2,1}$ normal error); like Lloyd (\S15) / $k$-planes (\S2) it
monotonically decreases the total $L^{2,1}$ distortion and converges. SDF segmentation clusters the
shape-diameter function (ray-cast inward) by a GMM + graph cut; sharp-edge/dihedral segmentation splits at
faces whose dihedral exceeds a threshold.
\end{proposition}
\begin{theorem}[Heat-method geodesics (Crane)]\label{thm:heat}
Geodesic distance $\phi$ from a source is recovered by: (1) solve $(\mathrm{Id}-t\,\Delta)u=\delta_{\text{src}}$
(one heat step); (2) normalize $X=-\nabla u/\Norm{\nabla u}$; (3) solve the Poisson equation
$\Delta\phi=\nabla\cdot X$. As $t\to0$, $\nabla u$ aligns with the geodesic direction (Varadhan
$u\sim e^{-\phi^2/4t}$), so $X$ is the unit geodesic gradient and step (3) integrates it to distance.
\end{theorem}
\begin{proof}[Sketch]
Varadhan's formula gives $\lim_{t\to0}-4t\log u_t=\phi^2$, hence $\nabla u_t\parallel-\nabla\phi$ with
$\Norm{\nabla\phi}=1$ (eikonal). Normalizing removes the magnitude error, and the least-squares integration of
a unit field is the Poisson problem of step (3). Discretized with the cotangent Laplacian.
\end{proof}

\section{Registration extensions}
\begin{proposition}[Weighted Kabsch = Horn]\label{prop:kabsch}
$\min_R\sum_i w_i\Norm{Rp_i'-q_i'}^2$ ($\sum w_i=1$, centred) is maximized by $R=V\,\mathrm{diag}(1,1,\det
VU^\top)\,U^\top$ where $H=\sum_i w_i p_i' q_i'^\top=U\Sigma V^\top$; this is the SVD form of Horn
(Thm~\ref{thm:horn}), with the $\det$ correction enforcing $R\in SO(3)$ (no reflection).
\end{proposition}
\begin{theorem}[CPD / Soft-ICP monotonicity]\label{thm:cpd}
Modeling the target as a Gaussian mixture over source points, EM alternates: E-step posteriors
$w_{ij}=\frac{\exp(-\Norm{q_j-Tp_i}^2/2\sigma^2)}{\sum_k\exp(-\Norm{q_j-Tp_k}^2/2\sigma^2)+c}$, M-step weighted
Kabsch (Prop~\ref{prop:kabsch}) + $\sigma^2$ update. Each EM round does not decrease the data
log-likelihood; hence Soft-ICP / CPD converges. Trimmed ICP instead keeps the best $h\%$ residuals (least
trimmed squares) and is monotone on the trimmed objective.
\end{theorem}
\begin{proof}
Standard EM: the E-step makes the variational bound tight, the M-step maximizes it; the likelihood is
non-decreasing and bounded, so it converges (Dempster--Laird--Rubin). ICP (\S11) is the $\sigma\to0$ hard-EM
limit.
\end{proof}
\begin{theorem}[Phase-correlation registration]\label{thm:phasecorr}
If two signals differ by a pure translation $t$, $f_2(x)=f_1(x-t)$, then their normalized cross-power
spectrum $\frac{\hat f_1\overline{\hat f_2}}{|\hat f_1\overline{\hat f_2}|}=e^{i\,k\cdot t}$, whose inverse
Fourier transform is $\delta(x-t)$. Hence the translation is the location of the inverse-FFT peak (exact for a
pure shift; robust to noise via the magnitude normalization). Extends to 2D/3D.
\end{theorem}
\begin{proof}
Fourier shift theorem: $\hat f_2(k)=\hat f_1(k)e^{-ik\cdot t}$, so $\hat f_1\overline{\hat f_2}=|\hat f_1|^2
e^{ik\cdot t}$; dividing by the magnitude leaves $e^{ik\cdot t}$, and $\mathcal F^{-1}[e^{ik\cdot t}]=\delta(x-t)$.
\end{proof}

\section{Scheduling, packing bounds, and graphs}
\begin{theorem}[Graham LPT bound]\label{thm:lpt}
Longest-processing-time-first list scheduling on $m$ identical machines yields makespan
$C_{\mathrm{LPT}}\le\big(\tfrac43-\tfrac{1}{3m}\big)C^\star$.
\end{theorem}
\begin{proof}[Sketch]
Consider the machine finishing last and the last job $j$ it runs; by LPT all jobs $\ge p_j$ were placed
first, and a counting argument on the load before $j$ plus $p_j\le C^\star/... $ gives the $4/3-1/(3m)$ ratio
(Graham 1969).
\end{proof}
\begin{proposition}[FFD bin packing]
First-Fit-Decreasing uses $\mathrm{FFD}\le\tfrac{11}{9}\mathrm{OPT}+\tfrac69$ bins (Johnson; Dosa tight
additive); the heightmap/tree 3D packers (\S7) are the geometric analogue.
\end{proposition}
\begin{proposition}[Welsh--Powell coloring]
Coloring greedily in non-increasing degree order uses at most $\max_i\min(d_i+1,\,i)\le\Delta+1$ colors
($d_i$ degree, $i$ the position); a proper coloring results since each vertex avoids its already-colored
neighbours.
\end{proposition}
\begin{theorem}[Kahn topological install order]\label{thm:kahn}
Repeatedly emitting and removing a zero-in-degree vertex of a directed graph $G$ produces a topological order
iff $G$ is acyclic; the order respects all support/precedence edges (a part is placed only after its
supports). If a non-empty residual has no source, $G$ has a cycle (no valid build order).
\end{theorem}
\begin{proof}
If $G$ is a DAG it has a source (else follow in-edges to form a cycle); removing it preserves acyclicity, so
by induction the process empties $G$ in an order where every edge $u\to v$ emits $u$ before $v$. Conversely a
cycle never reaches in-degree $0$, leaving a stuck non-empty residual.
\end{proof}
\begin{proposition}[NSGA-II]
Fast non-dominated sorting partitions the population into Pareto fronts in $O(MN^2)$ ($M$ objectives), and
crowding distance (the per-objective neighbour gap sum) breaks ties to preserve front diversity; selection by
$(\text{rank},\,-\text{crowding})$ drives the population to the Pareto front (\S15 dominance).
\end{proposition}

\section{Stereotomy and sphere projections}
\begin{definition}[Lines of curvature]
A line of curvature is an integral curve of a principal-direction field, i.e.\ a curve whose tangent is
everywhere an eigenvector of the shape operator $S=-dN$ (eigenvalues the principal curvatures
$\kappa_1,\kappa_2$). The two families form an orthogonal conjugate net.
\end{definition}
\begin{proposition}[Curvature-net voussoirs]
Bricks bounded by two families of curvature lines (Monge / Frézier stereotomy; pendentive and radial voussoir
cells) have faces that are near-planar developable strips, because along a line of curvature the surface
normal varies in a single plane (Rodrigues' formula $dN=-\kappa\,dx$), making the bed joints rule-aligned ---
the classical reason curvature-line tessellations are preferred for stone cutting.
\end{proposition}
\begin{theorem}[Equal-area (Lambert/Schmidt) projection]\label{thm:lambert}
The lower-hemisphere map $r=\sqrt2\,\sin(\theta/2)$ (colatitude $\theta$, azimuth $\varphi$) onto the disk is
area-preserving: it sends the spherical area element $\sin\theta\,d\theta\,d\varphi$ to $r\,dr\,d\varphi$ up to
the constant factor, so joint-pole \emph{density} on the Schmidt net is undistorted (unlike the equal-angle
Wulff net). Hence contoured pole densities give unbiased joint-set strengths (\S13).
\end{theorem}
\begin{proof}
With $r=\sqrt2\sin(\theta/2)$, $r\,dr=\sqrt2\sin(\theta/2)\cdot\sqrt2\cdot\tfrac12\cos(\theta/2)\,d\theta
=\sin(\theta/2)\cos(\theta/2)\,d\theta=\tfrac12\sin\theta\,d\theta$. Thus $r\,dr\,d\varphi=\tfrac12\sin\theta
\,d\theta\,d\varphi$, a constant multiple of the sphere's area element --- area preserved.
\end{proof}

\section{Robust statistics and supporting routines}
\begin{definition}[Tukey fence vs.\ $2\sigma$]
A value is a Tukey outlier iff outside $[Q_1-1.5\,\mathrm{IQR},\,Q_3+1.5\,\mathrm{IQR}]$ ($\mathrm{IQR}=Q_3-Q_1$),
a rank-based rule with $25\%$ breakdown, more robust than the mean-based $2\sigma$ clip of \S3 (which a single
gross outlier can distort). Both are used as pick/residual filters.
\end{definition}
\begin{proposition}[Supporting routines]
(i) Connected-components labelling via union--find runs in near-linear $O(n\,\alpha(n))$ and yields the
partition into maximal connected sets. (ii) BVH / triangle--AABB pruning is conservative: a query missing a
node's bounding box misses all its descendants (bounding-volume containment), so culling never discards a true
hit. (iii) Projective / position-based dynamics (Kangaroo) minimizes $\sum_c w_c\Norm{C_c(x)}^2$ by alternating
local constraint projections and a global solve, converging to a constraint-satisfying rest state --- a
relaxation of the contact KKT system of \S17.
\end{proposition}

\section*{Lean formalization notes}
\begin{itemize}[leftmargin=1.4em]
\item Polygons as \texttt{List} ($\R \times \R$); \texttt{area}, \texttt{ccw}, \texttt{reflex} as defs; clipping
as a structurally-recursive fold over edges. Lemma~\ref{lem:sh} (subset+area-monotone) is the key invariant.
\item Kriging reduces to SPD linear algebra (Mathlib \texttt{Matrix.PosDef}); Thm~\ref{thm:nugget} is a
Rayleigh/operator-norm bound; Thm~\ref{thm:kplanes} is Lloyd monotonicity over a finite assignment set.
\item Thm~\ref{thm:imaiiri} is a graph-path $\leftrightarrow$ approximation bijection + BFS optimality; the
cleanest to mechanize first.
\item Combinatorial/recursive targets: Thm~\ref{thm:nfp} (NFP non-overlap), Thm~\ref{thm:guillodp}
(guillotine tiling = Bellman recursion), Thm~\ref{thm:hm} (Hertel--Mehlhorn $4$-approx), and the 3D clip
Lemma~\ref{lem:clip3d} (subset invariant, cf.\ Lemma~\ref{lem:sh}).
\item Convex-analysis targets (Mathlib convex cones / KKT): Thm~\ref{thm:cra} (CRA stability = SOCP
feasibility), Thm~\ref{thm:settle} (drop-settle rest = KKT of constrained PE), Thm~\ref{thm:blocktheory}
(removability = joint/excavation cone emptiness).
\item Lie-group targets (Mathlib has $SO(3)$): Thm~\ref{thm:horn} (absolute orientation = top quaternion
eigenvector) and Thm~\ref{thm:kintsugi} (pose composition + global-pose invariance in $SE(3)/\mathrm{Sim}(3)$).
\item PDE/variational targets: Thm~\ref{thm:poisson} (Poisson reconstruction E--L), Thm~\ref{thm:stolt}
(Stolt remap, exact for constant $v$), and the CVT/GVF energy descents.
\item (\S20--27) Counting/induction targets: Thm~\ref{thm:lpt} (LPT $4/3$), FFD $11/9$, Welsh--Powell
$\Delta+1$, Thm~\ref{thm:potato} (convex skull $\ge$ greedy), Thm~\ref{thm:kahn} (Kahn on a DAG).
\item (\S20--27) Linear-algebra / Fourier targets: Thm~\ref{thm:qem} (QEM $=$ squared plane distances),
Prop~\ref{prop:kabsch} (weighted Kabsch SVD $=$ Horn), Thm~\ref{thm:phasecorr} (phase correlation $=$ shift
theorem), Prop~\ref{prop:power} (power cells convex), Prop~\ref{prop:pca} (PCA spectral).
\item (\S20--27) Analysis targets: Thm~\ref{thm:heat} (heat method via Varadhan), Thm~\ref{thm:cpd} (EM
monotonicity), Thm~\ref{thm:lambert} (Lambert area preservation $=$ a Jacobian identity, fully mechanizable).
\item Companion prose for older variants/back-ends in
\texttt{frahan-stonepack/research/MATH\_DERIVATIONS.md}; this file now supersedes it.
\end{itemize}
\end{document}
